\chapter{Nervous system diseases}
\index{Nervous system diseases}

\section*{Definition and Overview}
Nervous system diseases (or neurological disorders) encompass a vast array of pathological conditions affecting the brain, spinal cord, cranial nerves, peripheral nerves, nerve roots, autonomic nervous system, neuromuscular junctions, and skeletal muscles. These disorders disrupt the transmission of electrical and chemical signals within the body, leading to structural, electrical, or metabolic dysfunction. Understanding these diseases is critical, as they can cause profound physical, cognitive, and sensory impairment, demanding comprehensive clinical assessment and specialised management strategies. They range from acute events, such as stroke and traumatic brain injury, to chronic progressive conditions like Alzheimer's disease and multiple sclerosis.

\section*{Epidemiology}
Globally, nervous system diseases are the leading cause of disability-adjusted life years (DALYs) and the second leading cause of death, accounting for approximately 9 million deaths annually. In many countries, neurological disorders place an immense burden on the healthcare system. The prevalence of neurodegenerative diseases, such as Alzheimer's and other dementias, is increasing rapidly due to an ageing population. Chronic conditions like migraines affect approximately 15\% of the the population, while Parkinson's disease and multiple sclerosis show stable or rising age-standardised incidence rates, with multiple sclerosis having a higher prevalence in southern latitudes of some countries.

\section*{Aetiology and Pathophysiology}
The pathological mechanisms underlying nervous system diseases are diverse, reflecting the complexity of the nervous tissue and its vulnerability to various forms of injury.
\begin{itemize}
    \item \textbf{Vascular dysfunction:} Atherosclerosis, thrombosis, or arterial rupture causing ischaemic or haemorrhagic stroke, leading to localized tissue hypoxia.
    \item \textbf{Neurodegeneration:} Misfolding and aggregation of pathological proteins (e.g.\ beta-amyloid, tau, alpha-synuclein), leading to synaptic dysfunction, mitochondrial failure, and progressive neuronal death.
    \item \textbf{Autoimmune neuroinflammation:} T-cell and B-cell mediated attack on myelin or axonal components, as seen in multiple sclerosis or Guillain-Barré syndrome.
    \item \textbf{Infectious processes:} Pathogenic invasion of the meninges or brain parenchyma by bacteria (e.g.\ \textit{Neisseria meningitidis}), viruses (e.g.\ herpes simplex virus), or fungi.
    \item \textbf{Oncological processes:} Primary neuroepithelial tumours (e.g.\ glioblastoma), meningiomas, or metastatic lesions from other primary sites (e.g.\ lung, breast).
    \item \textbf{Toxic and metabolic insults:} Chronic alcohol toxicity, nutritional deficiencies (such as vitamin B12 deficiency), and heavy metal poisoning.
\end{itemize}

\section*{Clinical Features}
Clinical manifestations are dictated by the part of the nervous system affected (central, peripheral, or autonomic) and the rate of disease onset.
\begin{itemize}
    \item \textbf{Cognitive decline:} Loss of memory, language difficulties (aphasia), apraxia, agnosia, and executive dysfunction.
    \item \textbf{Motor impairment:} Weakness, spasticity, rigidity, tremor, chorea, ataxia, and loss of fine motor control.
    \item \textbf{Sensory disturbances:} Neuropathic pain, paraesthesia, hypoesthesia, and deficits in special senses such as vision, hearing, or smell.
    \item \textbf{Autonomic dysfunction:} Bladder and bowel incontinence, impotence, postural hypotension, and impaired thermoregulation.
    \item \textbf{Paroxysmal events:} Seizures, transient ischaemic attacks (TIAs), and acute headache episodes (migraines).
\end{itemize}

\section*{Differential Diagnosis}
Given the overlapping symptoms of neurological conditions, a rigorous differential diagnosis is essential to rule out non-neurological diseases.
\begin{itemize}
    \item \textbf{Psychogenic/functional neurological disorders:} Physical neurological symptoms (e.g.\ functional weakness or non-epileptic seizures) without organic pathology.
    \item \textbf{Systemic endocrine disorders:} Severe hypothyroidism or diabetes mellitus presenting with neuropathic symptoms or cognitive slowing.
    \item \textbf{Electrolyte imbalances:} Hyponatremia, hypokalaemia, or hypercalcaemia causing confusion, seizures, or muscle weakness.
    \item \textbf{Cardiovascular conditions:} Postural orthostatic tachycardia syndrome (POTS) or cardiac arrhythmias presenting with syncope or dizziness.
    \item \textbf{Rheumatological diseases:} Sjögren's syndrome or systemic lupus erythematosus presenting with peripheral neuropathy or neuropsychiatric symptoms.
\end{itemize}

\section*{When to Seek Medical Care}
Prompt medical consultation is necessary for any new, progressive, or unexplained neurological symptoms.
\hfill \break
\textbf{Contact your local emergency services for:}
\begin{itemize}
    \item Sudden weakness, numbness, or paralysis in the face, arm, or leg, especially on one side of the body.
    \item Sudden confusion, difficulty speaking, or difficulty understanding speech.
    \item Sudden visual loss, double vision, or severe dizziness and imbalance.
    \item A new, sudden, severe headache, or any seizure activity in someone without a known history of epilepsy.
\end{itemize}

\section*{Diagnosis and Investigations}
Diagnosis requires combining detailed history, physical examination, imaging, and neurophysiological testing.
\begin{itemize}
    \item \textbf{Clinical neurological examination:} Detailed assessment of cognitive state, cranial nerves, reflexes, motor strength, sensory pathways, and cerebellar function.
    \item \textbf{Magnetic resonance imaging (MRI):} Gold standard for visualising soft tissue structures, demyelinating lesions, ischaemic changes, and intracranial masses.
    \item \textbf{Computed tomography (CT):} Used in acute settings to rule out intracranial haemorrhage or large structural lesions rapidly.
    \item \textbf{Electrodiagnostic testing (EMG and NCS):} To evaluate peripheral nerve function and muscle health, distinguishing neurogenic from myogenic disorders.
    \item \textbf{Lumbar puncture (CSF analysis):} To look for inflammatory markers, oligoclonal bands, proteins, glucose, and pathogens in the cerebrospinal fluid.
\end{itemize}

\section*{Management}
Management plans must be comprehensive, multidisciplinary, and patient-centred, involving pharmacological, surgical, and supportive therapies.
\begin{itemize}
    \item \textbf{Pharmacotherapy:} Disease-modifying therapies (for multiple sclerosis), levodopa and dopamine agonists (for Parkinson's), anticonvulsants (for epilepsy), and triptans or CGRP inhibitors (for migraines).
    \item \textbf{Neurosurgical intervention:} Resection of tumours, clipping or coiling of aneurysms, spinal decompression, or implantation of deep brain stimulators.
    \item \textbf{Multidisciplinary rehabilitation:} Intensive physiotherapy, occupational therapy, and speech therapy to maximise functional capacity.
    \item \textbf{Lifestyle and nutritional support:} Implementing balanced diets, regular low-impact exercise, and cognitive training.
    \item \textbf{Psychological support:} Counselling, cognitive behavioural therapy, and support groups to manage the mental impact of living with a chronic disease.
\end{itemize}

\section*{Complications}
\begin{itemize}
    \item Progressive, irreversible physical disability and loss of mobility.
    \item Severe cognitive impairment and dementia, requiring full-time care.
    \item Neurogenic dysphagia, leading to aspiration pneumonia and malnutrition.
    \item Chronic neuropathic pain syndromes and refractory depression.
    \item Autonomic dysfunction-related complications, such as severe orthostatic hypotension or neurogenic bladder.
\end{itemize}

\section*{Prognosis and Outcomes}
Prognosis varies significantly depending on the specific disease. Some conditions, like migraines or well-controlled epilepsy, have a favourable prognosis with minimal long-term disability. In contrast, neurodegenerative diseases like motor neurone disease (MND) or Alzheimer's disease are progressive and currently incurable, with management focused on palliative care and quality of life. Early diagnosis and intervention remain the key factors in optimising long-term functional outcomes.

\section*{Prevention and Self-Care}
\begin{itemize}
    \item Manage cardiovascular risk factors (blood pressure, cholesterol, diabetes) to protect cerebrovascular health.
    \item Maintain a physically active lifestyle to support neuroplasticity and cardiovascular function.
    \item Protect the head from injury by using safety gear during sports and driving.
    \item Limit alcohol intake, avoid smoking, and avoid exposure to known neurotoxins.
    \item Adopt a nutrient-rich, balanced diet (such as the Mediterranean diet) and keep mentally stimulated through lifelong learning.
\end{itemize}

\section*{Sources and Further Reading}
The following organisations offer expert guidance and resources on nervous system diseases:
\begin{itemize}
    \item Brain Foundation. \textit{Support and research funding for brain and spinal cord disorders.} https://brainfoundation.org.au/
    \item MS Australia. \textit{National body for multiple sclerosis advocacy and information.} https://www.msaustralia.org.au/
    \item Dementia Australia. \textit{Peak body representing and supporting people living with dementia.} https://www.dementia.org.au/
    \item NHS. \textit{Comprehensive health directory on neurological conditions.} https://www.nhs.uk/
    \item Mayo Clinic. \textit{Clinical overview of nervous system diseases and therapies.} https://www.mayoclinic.org/
\end{itemize}
